% AUTO-GENERATED by paper/gen_data.py from docs/data/levels.json — do not edit.
\begin{longtable}{@{}c l p{8cm}@{}}
\toprule
Ур. & Запись & Насколько это много \\
\midrule
\endhead
a & $10^{36}$ & примерно как ундециллион (10\textasciicircum 36) — здесь начинается уровень a \\
b & $(10^{36})^{\star 1}$ & Обгоняет гугол и вообще любое физически осмысленное число. \\
c & $(10^{36})^{\star 2}$ & Уже больше гуголплекса — один только показатель степени астрономичен. \\
d & $(10^{36})^{\star 3}$ & У числа его цифр — больше цифр, чем атомов во Вселенной. \\
e & $(10^{36})^{\star 4}$ & Возьмите логарифм трижды подряд — всё ещё за пределами всей физики. \\
f & $(10^{36})^{\star 5}$ & Возведите гуголплекс в гуголплекс — вы даже не рядом. \\
g & $(10^{36})^{\star 6}$ & Больше шахматных партий Шеннона в степени числа позиций го. \\
h & $(10^{36})^{\star 7}$ & «Трудолюбивый бобёр» остановится раньше, чем напечатает первый показатель. \\
i & $(10^{36})^{\star 8}$ & Впишите число Авогадро в каждый планковский объём космоса — всё равно мало. \\
j & $(10^{36})^{\star 9}$ & Каждая цифра — вселенная, атомы каждой вселенной — цифры; всё равно мало. \\
k & $(10^{36})^{\star 10}$ & Само число его цифр уже за пределами всякого счёта. \\
l & $(10^{36})^{\star 11}$ & В глубине тетрации: самовозведения выше любого именованного числа. \\
m & $(10^{36})^{\star 12}$ & Функция Аккермана сочла бы это одним из первых, разминочных значений. \\
n & $(10^{36})^{\star 13}$ & Больше числа Скьюза — когда-то «крупнейшего числа в математике». \\
o & $(10^{36})^{\star 14}$ & Называйте новое число каждый планковский миг от Большого взрыва — не дойдёте. \\
p & $(10^{36})^{\star 15}$ & Гуголплексиан — лишь погрешность округления у его основания. \\
q & $(10^{36})^{\star 16}$ & Его логарифму нужен свой логарифм, а тому — ещё свой… \\
r & $(10^{36})^{\star 17}$ & Выше любой башни из десяток, что можно сложить атом за атомом. \\
s & $(10^{36})^{\star 18}$ & Показатель показателя показателя уже нефизичен. \\
t & $(10^{36})^{\star 19}$ & Дальше всех «-иллионов» и всех чисел, что люди когда-либо называли. \\
u & $(10^{36})^{\star 20}$ & Исчерпаете факториалы факториалов — и всё ещё не догоните. \\
v & $(10^{36})^{\star 21}$ & Двадцать раз «возведи в себя» — и вы едва добрались сюда. \\
w & $(10^{36})^{\star 22}$ & Башня степеней уже так высока, что десятки сливаются в одну линию. \\
x & $(10^{36})^{\star 23}$ & И всё же — отрезвляющая правда — бесконечно далеко от числа Грэма. \\
y & $(10^{36})^{\star 24}$ & Бесконечно далеко от TREE(3), рядом с которым число Грэма ничтожно. \\
z & $(10^{36})^{\star 25}$ & Последняя буква — алфавит кончился раньше, чем число. \\
$\ccoma$ & $(10^{36})^{\star 26}$ & Вершина — и всё же конечное число, которое в принципе можно записать. \\
\bottomrule
\end{longtable}
